\section{Space Complexity --- Padding}
%%% generated by the blueprint pipeline from TCSlib.Complexity.SpaceComplexity.Padding
This module proves the space form of the padding argument: if $L \in \mathrm{SPACE}(S)$,
then the language of members of $L$ padded by $C(|x|+1)^d$ ones lies in $\mathrm{SPACE}(T)$ for every
$T$ that is at least linear and dominates $S$ at the padded length. The proof builds a
logspace validator for padded words, a logspace bank answering queries about the first component,
and a simulation of the decider for $L$ on that component read as a virtual input.

\begin{definition}[Padded language]
\lean{Complexity.padLang}
\label{Complexity.padLang}
\leanok
\uses{Turing.pairEncode}
Let $L \subseteq \{0,1\}^*$ be a language and let $\mathrm{pad} : \mathbb{N} \to \mathbb{N}$.
The padded language $\mathrm{Pad}(L,\mathrm{pad})$ is the set of all strings
$\langle x, 1^{\mathrm{pad}(|x|)}\rangle$ with $x \in L$, where $\langle a,b\rangle$ denotes the self-delimiting pair encoding (every bit of $a$ doubled, then $01$, then $b$).
\end{definition}

\begin{theorem}[Membership in the padded language]
\lean{Complexity.mem_padLang_iff}
\label{Complexity.mem_padLang_iff}
\leanok
\uses{Complexity.padLang, Turing.pairEncode}
\difficulty{1}
Let $L \subseteq \{0,1\}^*$, let $\mathrm{pad} : \mathbb{N} \to \mathbb{N}$, and let
$y \in \{0,1\}^*$. Then $y$ belongs to the padded language $\mathrm{Pad}(L,\mathrm{pad})$ if and
only if there is some $x \in L$ with $y = \langle x, 1^{\mathrm{pad}(|x|)}\rangle$, where $\langle a,b\rangle$ denotes the self-delimiting pair encoding (every bit of $a$ doubled, then $01$, then $b$).
\end{theorem}

\begin{definition}[Padded words]
\lean{Complexity.epadValid}
\label{Complexity.epadValid}
\leanok
\uses{Turing.pairEncode}
Let $C, d \in \mathbb{N}$. A binary string $y$ is called $(C,d)$-padded if
$y = \langle a, 1^{C(|a|+1)^d}\rangle$ for some binary string $a$, where $\langle a,b\rangle$ denotes the self-delimiting pair encoding (every bit of $a$ doubled, then $01$, then $b$).
\end{definition}

\begin{definition}[Language of padded words]
\lean{Complexity.epadV}
\label{Complexity.epadV}
\leanok
\uses{Complexity.epadValid}
Let $C, d \in \mathbb{N}$. The language $V_{C,d}$ is the set of all $(C,d)$-padded binary
strings, that is, of all strings $\langle a, 1^{C(|a|+1)^d}\rangle$ with $a \in \{0,1\}^*$,
where $\langle a,b\rangle$ denotes the self-delimiting pair encoding (every bit of $a$ doubled, then $01$, then $b$).
\end{definition}

\begin{lemma}[Membership in the padded-word language]
\lean{Complexity.epad_mem_V}
\label{Complexity.epad_mem_V}
\leanok
\uses{Complexity.epadV, Complexity.epadValid}
\difficulty{1}
Let $C, d \in \mathbb{N}$ and $y \in \{0,1\}^*$. Then $y$ lies in the language $V_{C,d}$ of
$(C,d)$-padded words if and only if $y = \langle a, 1^{C(|a|+1)^d}\rangle$ for some binary
string $a$, where $\langle a,b\rangle$ denotes the self-delimiting pair encoding (every bit of $a$ doubled, then $01$, then $b$).
\end{lemma}

\begin{definition}[Virtual input of a padded word]
\lean{Complexity.epadVirt}
\label{Complexity.epadVirt}
\leanok
\uses{Complexity.epadValid, Complexity.pairFstD}
Let $C, d \in \mathbb{N}$. For a binary string $y$, the virtual input $\mathrm{virt}_{C,d}(y)$ is
the first component $a$ of $y$ when $y = \langle a, 1^{C(|a|+1)^d}\rangle$ is a
$(C,d)$-padded word, and the empty string otherwise; here $\langle a,b\rangle$ denotes the self-delimiting pair encoding (every bit of $a$ doubled, then $01$, then $b$).
\end{definition}

\begin{definition}[Query languages of the bank]
\lean{Complexity.epadQ}
\label{Complexity.epadQ}
\leanok
\uses{Complexity.LogProg.bitsVal, Complexity.epadValid, Complexity.pairFstD, Turing.pairDecode}
Let $C, d \in \mathbb{N}$ and let $\mu \in \{0,1\}$ be a mode. The language $Q_{C,d,\mu}$
consists of the strings $z = \langle y, w\rangle$ (with $\langle\cdot,\cdot\rangle$ the
self-delimiting pair encoding) such that $w$ is empty or ends in $1$, $|w|$ is at most the length
of the binary representation of $|y|$, and $y = \langle a, 1^{C(|a|+1)^d}\rangle$ is
$(C,d)$-padded, and moreover: in mode $\mu = 0$, $\mathrm{val}(w) < |a|$; in mode $\mu = 1$, the
letter of $a$ at position $\mathrm{val}(w)$ is $1$ (letters beyond the end count as $0$). Here
$\mathrm{val}(w)$ is the value of $w$ read as a little-endian binary numeral.
\end{definition}

\begin{lemma}[Virtual input of a padded word is its first component]
\lean{Complexity.epadVirt_of_valid}
\label{Complexity.epadVirt_of_valid}
\leanok
\uses{Complexity.epadValid, Complexity.epadVirt, Complexity.pairFstD}
\difficulty{1}
Let $C, d \in \mathbb{N}$ and let $y$ be a $(C,d)$-padded binary string, i.e.
$y = \langle a, 1^{C(|a|+1)^d}\rangle$ for some $a$ in the self-delimiting pair encoding. Then
the virtual input of $y$ (the first component of $y$ if $y$ is $(C,d)$-padded, the empty string
otherwise) equals the first component $a$ of $y$.
\end{lemma}

\begin{lemma}[Virtual input of an unpadded word is empty]
\lean{Complexity.epadVirt_of_not_valid}
\label{Complexity.epadVirt_of_not_valid}
\leanok
\uses{Complexity.epadValid, Complexity.epadVirt}
\difficulty{1}
Let $C, d \in \mathbb{N}$ and let $y$ be a binary string that is not of the form
$\langle a, 1^{C(|a|+1)^d}\rangle$ in the self-delimiting pair encoding. Then the virtual input
of $y$ (the first component of $y$ if $y$ has this form, the empty string otherwise) is the
empty string.
\end{lemma}

\begin{lemma}[Virtual input is no longer than the input]
\lean{Complexity.epadVirt_length_le}
\label{Complexity.epadVirt_length_le}
\leanok
\uses{Complexity.epadValid, Complexity.epadVirt, Complexity.epadVirt_of_not_valid, Complexity.epadVirt_of_valid, Complexity.length_pairFstD_le}
\difficulty{3}
Let $C, d \in \mathbb{N}$ and $y \in \{0,1\}^*$. The virtual input of $y$ (its first component
$a$ if $y = \langle a, 1^{C(|a|+1)^d}\rangle$ in the self-delimiting pair encoding, and the
empty string otherwise) has length at most $|y|$.
\end{lemma}

\begin{lemma}[Membership in a bank language for a decoded pair]
\lean{Complexity.epad_mem_Q_iff}
\label{Complexity.epad_mem_Q_iff}
\leanok
\uses{Complexity.LogProg.bitsVal, Complexity.epadQ, Complexity.epadValid, Complexity.pairFstD, Turing.pairDecode}
\difficulty{3}
Let $C, d \in \mathbb{N}$, let $\mu \in \{0,1\}$, and let $z$ be a binary string that decodes as
the self-delimiting pair $z = \langle y, w\rangle$. Then $z \in Q_{C,d,\mu}$ if and only if $w$ is
empty or ends in $1$, $|w|$ is at most the bit length of $|y|$, $y$ is $(C,d)$-padded (that is,
$y = \langle a, 1^{C(|a|+1)^d}\rangle$ for some $a$), and the mode condition holds, where $a$ is the first component of $y$: if $\mu = 1$ the letter of $a$ at position $\mathrm{val}(w)$
is $1$, and if $\mu = 0$ then $\mathrm{val}(w) < |a|$, where $\mathrm{val}(w)$ is the
little-endian binary value of $w$.
\end{lemma}

\begin{lemma}[Non-pairs are not in a bank language]
\lean{Complexity.epad_not_mem_Q}
\label{Complexity.epad_not_mem_Q}
\leanok
\uses{Complexity.epadQ, Turing.pairDecode}
\difficulty{2}
Let $C, d \in \mathbb{N}$ and $\mu \in \{0,1\}$. If a binary string $z$ does not decode as a
self-delimiting pair $\langle y, w\rangle$, then $z$ does not belong to the bank language
$Q_{C,d,\mu}$, whose members are all of this pair form.
\end{lemma}

\begin{lemma}[Bank answers on genuine queries]
\lean{Complexity.epad_query_mem_iff}
\label{Complexity.epad_query_mem_iff}
\leanok
\uses{Complexity.LogProg.bitsVal_bits, Complexity.LogProg.canon_bits, Complexity.LogProg.drop_canon, Complexity.LogProg.length_bits_mono, Complexity.epadQ, Complexity.epadValid, Complexity.epad_mem_Q_iff, Complexity.pairFstD, Turing.pairDecode_pairEncode, Turing.pairEncode}
\difficulty{3}
Let $C, d, i \in \mathbb{N}$, $\mu \in \{0,1\}$ and $y \in \{0,1\}^*$ with $i \le |y|$, and write
$\mathrm{bin}(i)$ for the binary representation of $i$. Then
$\langle y, \mathrm{bin}(i)\rangle \in Q_{C,d,\mu}$ if and only if $y$ is $(C,d)$-padded, i.e.
$y = \langle a, 1^{C(|a|+1)^d}\rangle$ for some $a$, and, for its first component $a$: in mode
$\mu = 1$ the letter $a_i$ is $1$ (letters beyond the end count as $0$), and in mode $\mu = 0$
we have $i < |a|$.
\end{lemma}

\begin{definition}[Instructions with unconditional step preconditions]
\lean{Complexity.LogProg.epadOk}
\label{Complexity.LogProg.epadOk}
\leanok
\uses{Complexity.LogProg.Ins}
Let $\iota$ be an instruction of an abstract register machine. The predicate
$\mathrm{Ok}(\iota)$ holds exactly when either $\iota$ is an equality test between two distinct
registers, or $\iota$ is a decider call whose argument registers are pairwise distinct, or
$\iota$ is any other instruction except the three tests that compare a register's binary
representation with a component of the input (for which it fails).
\end{definition}

\begin{lemma}[Step preconditions from a well-formed instruction]
\lean{Complexity.LogProg.epad_preS}
\label{Complexity.LogProg.epad_preS}
\leanok
\uses{Complexity.LogProg.ARM, Complexity.LogProg.Ins, Complexity.LogProg.PreS, Complexity.LogProg.epadOk, Complexity.LogProg.preS_of_ins}
\difficulty{2}
Let $A$ be an abstract register machine, $x$ an input, and $\ell$ a label whose instruction is an
equality test of two distinct registers, a decider call with pairwise distinct argument
registers, or any instruction other than the three tests comparing a register with a component
of the input. Then every configuration of $A$ at label $\ell$, with arbitrary register contents
and answer, satisfies the syntactic step preconditions of $A$ on $x$ (equality tests compare
distinct registers, decider calls use pairwise distinct argument registers, and input-comparison
tests are applied only to well-formed inputs).
\end{lemma}

\begin{theorem}[Correctness of the ones-scanning loop]
\lean{Complexity.LogProg.epad_ones_loop}
\label{Complexity.LogProg.epad_ones_loop}
\leanok
\uses{Complexity.LogProg.ARM, Complexity.LogProg.AReach, Complexity.LogProg.AReach.single, Complexity.LogProg.AReach.step, Complexity.LogProg.AReach.trans, Complexity.LogProg.PreS, Complexity.LogProg.arm_chain, Complexity.LogProg.astep_call, Complexity.LogProg.astep_inc, Complexity.LogProg.astep_jeq, Complexity.LogProg.epad_preS, Complexity.LogProg.register_update_le, Complexity.LogProg.vword_whole₁, Turing.pairEncode}
\difficulty{5}
Let $A$ be an abstract register machine with an oracle answering its decider calls, run on input
$x$. Let $p, \ell, c$ be pairwise distinct registers, $j$ a decider, $s \ge 1$ and $B$ natural
numbers, and suppose that label $\mathsf{test}$ jumps to $\mathsf{exit}$ if $p = \ell$ and to
$\mathsf{read}$ otherwise, $\mathsf{read}$ calls decider $j$ on $\langle x, \mathrm{bin}(p)\rangle$
and goes to $\mathsf{inc}$ on $1$ and to $\mathsf{fail}$ on $0$, $\mathsf{inc}$ increments $c$ and
goes to $\mathsf{adv}_0$, each $\mathsf{adv}_i$ with $i < s$ increments $p$ and goes to
$\mathsf{adv}_{i+1}$, and $\mathsf{adv}_s = \mathsf{test}$. Then for every string $w$ and register
contents $v$ with all registers at most $B$, $v(\ell) = v(p) + s|w|$, $v(c) + |w| \le B$, and
oracle answers $w_i$ for decider $j$ on $\langle x, \mathrm{bin}(v(p) + s i)\rangle$ for all
$i < |w|$: if every letter of $w$ is $1$, the machine runs from $\mathsf{test}$ with $v$ to
$\mathsf{exit}$ with $p$ set to $v(\ell)$ and $c$ set to $v(c) + |w|$; and if some letter of $w$ is
$0$, it runs from $\mathsf{test}$ with $v$ to $\mathsf{fail}$ with some register contents that are
at most $B$ and agree with $v$ outside $p$ and $c$. In both cases every configuration along the run satisfies the syntactic step preconditions (equality tests compare distinct registers and decider calls use pairwise distinct argument registers) and has all register values at most $B$.
\end{theorem}

\begin{theorem}[Correctness of the multiplication loop]
\lean{Complexity.LogProg.epad_mul_loop}
\label{Complexity.LogProg.epad_mul_loop}
\leanok
\uses{Complexity.LogProg.ARM, Complexity.LogProg.AReach, Complexity.LogProg.AReach.single, Complexity.LogProg.AReach.step, Complexity.LogProg.AReach.trans, Complexity.LogProg.PreS, Complexity.LogProg.astep_dec, Complexity.LogProg.astep_jz, Complexity.LogProg.epad_preS, Complexity.LogProg.registerCopy_loop, Complexity.LogProg.registerScale_loop, Complexity.LogProg.register_update_le}
\difficulty{5}
Let $A$ be an abstract register machine with an oracle, run on input $x$, let $k, t, p, q$ be
pairwise distinct registers and $B \in \mathbb{N}$, and suppose its labels are programmed as
follows: $\mathsf{ml}$ jumps to $\mathsf{exit}$ if $k = 0$ and to $\mathsf{md}$ otherwise;
$\mathsf{md}$ decrements $k$ and goes to $\mathsf{addt}$; $\mathsf{addt}$ jumps to $\mathsf{sct}$ if
$t = p$ and to $\mathsf{addb}$ otherwise; $\mathsf{addb}$ increments $t$ and goes to
$\mathsf{addt}$; $\mathsf{sct}$ jumps to $\mathsf{ml}$ if $t = 0$ and to $\mathsf{scd}$ otherwise;
$\mathsf{scd}$ decrements $t$ and goes to $\mathsf{sci}$; $\mathsf{sci}$ increments $q$ and goes to
$\mathsf{sct}$. Then for every $K \in \mathbb{N}$ and register contents $v$ with $v(k) = K$,
$v(t) = 0$, all registers at most $B$ and $v(q) + K\,v(p) \le B$, the machine runs from
$\mathsf{ml}$ with $v$ to $\mathsf{exit}$ with $k$ set to $0$, $q$ set to $v(q) + K\,v(p)$ and all
other registers unchanged, and every configuration along the run satisfies the syntactic step preconditions (equality tests compare distinct registers and decider calls use pairwise distinct argument registers) and has all register values at most $B$.
\end{theorem}

\begin{theorem}[Correctness of the powering loop]
\lean{Complexity.LogProg.epad_pow_loop}
\label{Complexity.LogProg.epad_pow_loop}
\leanok
\uses{Complexity.LogProg.ARM, Complexity.LogProg.AReach, Complexity.LogProg.AReach.single, Complexity.LogProg.AReach.step, Complexity.LogProg.AReach.trans, Complexity.LogProg.PreS, Complexity.LogProg.astep_clr, Complexity.LogProg.astep_dec, Complexity.LogProg.astep_jz, Complexity.LogProg.epad_mul_loop, Complexity.LogProg.epad_preS, Complexity.LogProg.registerCopy_loop, Complexity.LogProg.register_update_le}
\difficulty{6}
Let $A$ be an abstract register machine with an oracle, run on input $x$, let $B \in \mathbb{N}$,
and let $n, e, k, t, p, q$ be registers such that $k, t, p, q$ are pairwise distinct, $e$ differs
from $n, k, t, p, q$, and $n$ differs from $k, p, q$. Suppose its labels are programmed as follows:
$\mathsf{pw}$ jumps to $\mathsf{exit}$ if $e = 0$ and to $\mathsf{pd}$ otherwise; $\mathsf{pd}$
decrements $e$ and goes to $\mathsf{mktest}$; $\mathsf{mktest}$ jumps to $\mathsf{ml}$ if $k = n$
and to $\mathsf{mkbump}$ otherwise; $\mathsf{mkbump}$ increments $k$ and goes to $\mathsf{mktest}$;
the multiplication block $\mathsf{ml}, \mathsf{md}, \mathsf{addt}, \mathsf{addb}, \mathsf{sct},
\mathsf{scd}, \mathsf{sci}$ (which adds $k \cdot p$ to $q$, using $t$ as scratch) exits to
$\mathsf{mv}$; $\mathsf{mv}$ clears $p$ and goes to $\mathsf{mvtest}$; $\mathsf{mvtest}$ jumps to
$\mathsf{mvc}$ if $p = q$ and to $\mathsf{mvbump}$ otherwise; $\mathsf{mvbump}$ increments $p$ and
goes to $\mathsf{mvtest}$; and $\mathsf{mvc}$ clears $q$ and goes to $\mathsf{pw}$. Then for every
$E \in \mathbb{N}$ and register contents $v$ with $v(e) = E$, $v(q) = v(k) = v(t) = 0$,
$v(n) \ge 1$, all registers at most $B$ and $v(p)\,v(n)^E \le B$, the machine runs from
$\mathsf{pw}$ with $v$ to $\mathsf{exit}$ with $e$ set to $0$, $p$ set to $v(p)\,v(n)^E$ and all
other registers unchanged, and every configuration along the run satisfies the syntactic step preconditions (equality tests compare distinct registers and decider calls use pairwise distinct argument registers) and has all register values at most $B$.
\end{theorem}

\begin{definition}[Labels of the validity block]
\lean{Complexity.EPadVL}
\label{Complexity.EPadVL}
\leanok
Let $C, d \in \mathbb{N}$. The label set of the validity block is the finite set consisting of the
pair-scan labels $\mathsf{scan}$, $\mathsf{first}$, $\mathsf{test}_b$, $\mathsf{read}_b$ and
$\mathsf{adv}_{k,\sigma}$ ($b, \sigma \in \{0,1\}$, $k < 4$), the copy-and-halve labels
$\mathsf{cptest}, \mathsf{cpbump}, \mathsf{half}_1, \mathsf{half}_2$, the ones-scan labels
$\mathsf{ostest}, \mathsf{osread}, \mathsf{osinc}, \mathsf{osadv}_\sigma$, the constant-loading
labels $\mathsf{cinc}_i$ ($i < C$) and $\mathsf{dinc}_i$ ($i < d$), fifteen labels of the power
loop, and the comparison label $\mathsf{cmp}$.
\end{definition}

\begin{definition}[Enumeration of the validity-block labels]
\lean{Complexity.epadVLAll}
\label{Complexity.epadVLAll}
\leanok
\uses{Complexity.EPadVL}
Let $C, d \in \mathbb{N}$. The enumeration of the validity block's labels is the explicit finite
list containing each pair-scan, copy-and-halve and ones-scan label, the labels
$\mathsf{cinc}_i$ for $i < C$, the labels $\mathsf{dinc}_i$ for $i < d$, all power-loop labels and
the comparison label.
\end{definition}

\begin{definition}[Finiteness of the validity-block labels]
\lean{Complexity.epadVLFintype}
\label{Complexity.epadVLFintype}
\leanok
\uses{Complexity.EPadVL, Complexity.epadVLAll}
Let $C, d \in \mathbb{N}$. The set of labels of the validity block is given the structure of a
finite set, with every label appearing in the explicit enumeration of the validity block's
labels.
\end{definition}

\begin{definition}[Stride of the validity block]
\lean{Complexity.epadS}
\label{Complexity.epadS}
\leanok
For a flag $\sigma \in \{0,1\}$, the stride $s(\sigma)$ is defined to be $2$ if $\sigma = 1$ (the
doubled first field of a query) and $1$ if $\sigma = 0$ (the input itself).
\end{definition}

\begin{definition}[Targets of the pair-scan advance chains]
\lean{Complexity.epadAdvEnd}
\label{Complexity.epadAdvEnd}
\leanok
\uses{Complexity.EPadVL}
Let $C, d \in \mathbb{N}$ and let $\iota$ map the validity-block labels into a label set
$\Lambda$. For $k \in \{0,1,2,3\}$, the end target of advance chain $k$ is the image under $\iota$
of $\mathsf{test}_0$ if $k = 0$, of $\mathsf{test}_1$ if $k = 1$, of $\mathsf{scan}$ if $k = 2$, and
of $\mathsf{cptest}$ if $k = 3$.
\end{definition}

\begin{definition}[Pair-scan advance chains]
\lean{Complexity.epadAdv}
\label{Complexity.epadAdv}
\leanok
\uses{Complexity.EPadVL, Complexity.epadAdvEnd}
Let $C, d \in \mathbb{N}$, let $\iota$ map the validity-block labels into a label set $\Lambda$,
let $\sigma \in \{0,1\}$ be the stride flag and $k \in \{0,1,2,3\}$. Step $i$ of advance chain $k$
is the label $\iota(\mathsf{adv}_{k,0})$ if $i = 0$, the label $\iota(\mathsf{adv}_{k,1})$ if
$i = 1$ and $\sigma = 1$, and otherwise the end target of chain $k$, namely the image under
$\iota$ of $\mathsf{test}_0$, $\mathsf{test}_1$, $\mathsf{scan}$ or $\mathsf{cptest}$ for
$k = 0, 1, 2, 3$ respectively.
\end{definition}

\begin{definition}[Ones-scan advance chain]
\lean{Complexity.epadOsAdv}
\label{Complexity.epadOsAdv}
\leanok
\uses{Complexity.EPadVL}
Let $C, d \in \mathbb{N}$, let $\iota$ map the validity-block labels into a label set $\Lambda$,
and let $\sigma \in \{0,1\}$ be the stride flag. Step $i$ of the ones-scan advance chain is
$\iota(\mathsf{osadv}_0)$ if $i = 0$, $\iota(\mathsf{osadv}_1)$ if $i = 1$ and $\sigma = 1$, and
$\iota(\mathsf{ostest})$ otherwise.
\end{definition}

\begin{definition}[Chain loading the exponent]
\lean{Complexity.epadDI}
\label{Complexity.epadDI}
\leanok
\uses{Complexity.EPadVL}
Let $C, d \in \mathbb{N}$ and let $\iota$ map the validity-block labels into a label set
$\Lambda$. Step $j$ of the chain loading $d$ is $\iota(\mathsf{dinc}_j)$ if $j < d$ and
$\iota(\mathsf{pw})$, the entry of the power loop, otherwise.
\end{definition}

\begin{definition}[Chain loading the coefficient]
\lean{Complexity.epadCI}
\label{Complexity.epadCI}
\leanok
\uses{Complexity.EPadVL, Complexity.epadDI}
Let $C, d \in \mathbb{N}$ and let $\iota$ map the validity-block labels into a label set
$\Lambda$. Step $j$ of the chain loading $C$ is $\iota(\mathsf{cinc}_j)$ if $j < C$, and otherwise
the first step of the chain loading $d$, namely $\iota(\mathsf{dinc}_0)$ if $d > 0$ and
$\iota(\mathsf{pw})$ if $d = 0$.
\end{definition}

\begin{definition}[Code of the validity block]
\lean{Complexity.epadValCode}
\label{Complexity.epadValCode}
\leanok
\uses{Complexity.EPadVL, Complexity.LogProg.Ins, Complexity.epadAdvEnd, Complexity.epadCI, Complexity.epadDI}
Let $\sigma \in \{0,1\}$ be a stride flag with stride $s = s(\sigma) \in \{1,2\}$, let
$C, d \in \mathbb{N}$, let $\iota$ map the validity-block labels into a label set $\Lambda$, and
let $\mathsf{yes}, \mathsf{no} \in \Lambda$. The validity block assigns to each validity-block
label an instruction of an abstract register machine with $18$ registers and $2$ deciders: using
register $0$ as pointer, register $1$ as limit and decider $0$ to read the letter at the
pointer, it parses the scanned word as a pair $\langle a,b\rangle$ (going to $\mathsf{no}$ if it
is malformed), advancing the pointer by $s$ per letter; it copies the pointer into register $2$
and halves it (once if $s = 1$, twice if $s = 2$) to obtain $|a|+1$; it scans $b$, going to $\mathsf{no}$ on a
$0$ and counting its ones in register $3$; it loads $C$ into register $4$ and $d$ into register
$8$ by chains of increments; it runs the powering loop (scratch registers $5$, $6$, $7$) to make
register $4$ equal to $C(|a|+1)^d$; and it finally goes to $\mathsf{yes}$ if registers $3$ and $4$
are equal and to $\mathsf{no}$ otherwise.
\end{definition}

\begin{lemma}[Padding test for a decoded pair]
\lean{Complexity.epad_valid_iff_of_decode}
\label{Complexity.epad_valid_iff_of_decode}
\leanok
\uses{Complexity.epadValid, Turing.eq_pairEncode_of_pairDecode, Turing.pairDecode, Turing.pairDecode_pairEncode}
\difficulty{3}
Let $C, d \in \mathbb{N}$ and let $y$ be a binary string that decodes as the self-delimiting pair
$y = \langle a, b\rangle$. Then $y$ is $(C,d)$-padded, i.e. of the form
$\langle a', 1^{C(|a'|+1)^d}\rangle$ for some $a'$, if and only if $b = 1^{C(|a|+1)^d}$.
\end{lemma}

\begin{lemma}[Non-pairs are not padded]
\lean{Complexity.epad_not_valid_of_decode}
\label{Complexity.epad_not_valid_of_decode}
\leanok
\uses{Complexity.epadValid, Turing.pairDecode, Turing.pairDecode_pairEncode}
\difficulty{2}
Let $C, d \in \mathbb{N}$. If a binary string $y$ does not decode as a self-delimiting pair
$\langle a, b\rangle$, then $y$ is not $(C,d)$-padded, i.e. not of the form
$\langle a, 1^{C(|a|+1)^d}\rangle$ for any $a$.
\end{lemma}

\begin{lemma}[Characterisation of a run of ones]
\lean{Complexity.epad_eq_replicate_iff}
\label{Complexity.epad_eq_replicate_iff}
\leanok
\difficulty{3}
Let $b$ be a binary string and $n \in \mathbb{N}$. Then $b = 1^n$ if and only if every letter
$b_i$ with $i < |b|$ is $1$ and $|b| = n$.
\end{lemma}

\begin{theorem}[Correctness of the validity block]
\lean{Complexity.epad_val_run}
\label{Complexity.epad_val_run}
\leanok
\uses{Complexity.EPadVL, Complexity.LogProg.ARM, Complexity.LogProg.AReach, Complexity.LogProg.AReach.refl, Complexity.LogProg.AReach.single, Complexity.LogProg.AReach.step, Complexity.LogProg.AReach.trans, Complexity.LogProg.PreS, Complexity.LogProg.alignedPair_loop, Complexity.LogProg.arm_chain, Complexity.LogProg.astep_half, Complexity.LogProg.astep_jeq, Complexity.LogProg.epad_ones_loop, Complexity.LogProg.epad_pow_loop, Complexity.LogProg.epad_preS, Complexity.LogProg.registerCopy_loop, Complexity.LogProg.register_update_le, Complexity.epadAdv, Complexity.epadCI, Complexity.epadDI, Complexity.epadOsAdv, Complexity.epadS, Complexity.epadValCode, Complexity.epadValid, Complexity.epad_eq_replicate_iff, Complexity.epad_not_valid_of_decode, Complexity.epad_valid_iff_of_decode, Complexity.pairEncode_getD_snd, Complexity.pairFstD, Complexity.pairFstD_pairEncode, Turing.eq_pairEncode_of_pairDecode, Turing.length_pairEncode, Turing.pairDecode, Turing.pairEncode}
\difficulty{7}
Let $A$ be an abstract register machine with $18$ registers and $2$ deciders and an oracle
answering its calls, run on input $x$, let $C, d \in \mathbb{N}$, let $\sigma \in \{0,1\}$ with
stride $s = s(\sigma)$ ($1$ if $\sigma = 0$, $2$ if $\sigma = 1$), and suppose that $A$, on the
image of an embedding $\iota$ of the validity-block labels, carries the validity-block code for
$\sigma, C, d$ with exits $\mathsf{yes}, \mathsf{no}$. Let $y$ be a binary string such that decider
$0$ answers $y_j$ on $\langle x, \mathrm{bin}(s j)\rangle$ for all $j < |y|$, let
$B \ge C(|y|+1)^d + d + 2|y| + 2$, and let $v$ be register contents with all registers at most
$B$, $v(0) = 0$, $v(1) = s|y|$ and $v(r) = 0$ for $2 \le r \le 8$. Then there are register contents
$v'$, agreeing with $v$ on registers $9, \dots, 17$ and with all registers at most $B$, such that
the machine runs from $\iota(\mathsf{scan})$ with $v$ to $\mathsf{yes}$ with $v'$ if
$y = \langle a, 1^{C(|a|+1)^d}\rangle$ for some $a$ (in which case $v'(2) = |a| + 1$) and to
$\mathsf{no}$ with $v'$ otherwise, where every configuration along the run satisfies the syntactic step preconditions (equality tests compare distinct registers and decider calls use pairwise distinct argument registers) and has all register values at most $B$.
\end{theorem}

\begin{definition}[Extra labels of the validator]
\lean{Complexity.EPadVX}
\label{Complexity.EPadVX}
\leanok
The extra labels of the validator are the elements of the four-element set
$\{\mathsf{lscan}, \mathsf{linc}, \mathsf{yes}, \mathsf{no}\}$: the two labels of the input-length
count and the two answering labels.
\end{definition}

\begin{definition}[Finiteness of the validator's extra labels]
\lean{Complexity.epadVXFintype}
\label{Complexity.epadVXFintype}
\leanok
\uses{Complexity.EPadVX}
The set $\{\mathsf{lscan}, \mathsf{linc}, \mathsf{yes}, \mathsf{no}\}$ of extra labels of the
validator is given the structure of a finite set, enumerated by the list of its four
elements.
\end{definition}

\begin{definition}[Validator for padded words]
\lean{Complexity.epadVA}
\label{Complexity.epadVA}
\leanok
\uses{Complexity.EPadVL, Complexity.EPadVX, Complexity.LogProg.ARM, Complexity.epadValCode}
Let $C, d \in \mathbb{N}$. The validator is the abstract register machine with $18$ registers and
$2$ deciders whose labels are the validity-block labels together with
$\mathsf{lscan}, \mathsf{linc}, \mathsf{yes}, \mathsf{no}$: label $\mathsf{lscan}$ calls decider
$1$ on $\langle x, \mathrm{bin}(r_1)\rangle$ and goes to $\mathsf{linc}$ on $1$ and to the block's
$\mathsf{scan}$ on $0$; $\mathsf{linc}$ increments register $1$ and returns to $\mathsf{lscan}$; the
validity-block labels carry the validity-block code at stride $1$ with exits $\mathsf{yes}$ and
$\mathsf{no}$; and $\mathsf{yes}$, $\mathsf{no}$ halt with answers $1$ and $0$.
\end{definition}

\begin{definition}[Input deciders of the register machines]
\lean{Complexity.epadIn}
\label{Complexity.epadIn}
\leanok
\uses{Complexity.inputBitLang, Complexity.inputLenLang}
The pair of input languages consists, at index $0$, of the bit language
$\{\langle u, \mathrm{bin}(i)\rangle : u_i = 1\}$ (letters beyond the end of $u$ count as $0$)
and, at index $1$, of the length language $\{\langle u, \mathrm{bin}(i)\rangle : i < |u|\}$, where
$\langle\cdot,\cdot\rangle$ is the self-delimiting pair encoding and $\mathrm{bin}(i)$ is the
binary representation of $i$.
\end{definition}

\begin{lemma}[Input deciders are logspace]
\lean{Complexity.epadIn_mem}
\label{Complexity.epadIn_mem}
\leanok
\uses{Complexity.LOGSPACE, Complexity.epadIn, Complexity.inputBitLang, Complexity.inputBitLang_mem, Complexity.inputLenLang, Complexity.inputLenLang_mem}
\difficulty{2}
For each $j \in \{0,1\}$, the $j$-th input language (for $j = 0$ the bit language
$\{\langle u, \mathrm{bin}(i)\rangle : u_i = 1\}$, for $j = 1$ the length language
$\{\langle u, \mathrm{bin}(i)\rangle : i < |u|\}$) belongs to $\mathsf{l}$, the class of languages
decidable in logarithmic space.
\end{lemma}

\begin{definition}[Oracle of the input deciders]
\lean{Complexity.epadOr}
\label{Complexity.epadOr}
\leanok
\uses{Complexity.epadIn, Turing.MultiTapeTM.indicator}
The oracle of the register machines sends a decider index $j \in \{0,1\}$ and a binary string $V$
to $1$ if $V$ belongs to the $j$-th input language (the bit language for $j = 0$, the length
language for $j = 1$) and to $0$ otherwise.
\end{definition}

\begin{lemma}[Answers of the input deciders]
\lean{Complexity.epadIn_answer}
\label{Complexity.epadIn_answer}
\leanok
\uses{Complexity.epadOr, Complexity.indicator_inputBitLang, Complexity.indicator_inputLenLang, Complexity.inputBitLang, Complexity.inputLenLang, Turing.MultiTapeTM.indicator, Turing.pairEncode}
\difficulty{1}
Let $y$ be a binary string and $p \in \mathbb{N}$, and let $\mathrm{bin}(p)$ be the binary
representation of $p$. On the query $\langle y, \mathrm{bin}(p)\rangle$, the oracle of decider $0$
(the bit language) returns the letter $y_p$, taken to be $0$ when $p \ge |y|$, and the oracle of
decider $1$ (the length language) returns $1$ exactly when $p < |y|$.
\end{lemma}

\begin{lemma}[Polynomial register bound]
\lean{Complexity.epad_bound}
\label{Complexity.epad_bound}
\leanok
\difficulty{4}
For all natural numbers $C$, $d$ and $n$,
\[ C(n+1)^d + d + 4n + 4 \le (C + d + 8)(n+1)^{d+1}. \]
\end{lemma}

\begin{theorem}[Padded words are decidable in logspace]
\lean{Complexity.epad_V_mem}
\label{Complexity.epad_V_mem}
\leanok
\uses{Complexity.LOGSPACE, Complexity.LogProg.AHalt.ret, Complexity.LogProg.AReach.halt, Complexity.LogProg.arm_decides_poly, Complexity.LogProg.epad_preS, Complexity.LogProg.indicator_eq_decide, Complexity.LogProg.inputLength_loop, Complexity.LogProg.register_update_le, Complexity.epadIn, Complexity.epadIn_answer, Complexity.epadIn_mem, Complexity.epadOr, Complexity.epadS, Complexity.epadV, Complexity.epadVA, Complexity.epadVLFintype, Complexity.epadVXFintype, Complexity.epadValid, Complexity.epad_bound, Complexity.epad_mem_V, Complexity.epad_val_run, Turing.MultiTapeTM.indicator}
\difficulty{5}
Let $C, d \in \mathbb{N}$. The language of $(C,d)$-padded words, i.e. of all strings
$\langle a, 1^{C(|a|+1)^d}\rangle$ with $a \in \{0,1\}^*$ in the self-delimiting pair encoding,
belongs to $\mathsf{l}$, the class of languages decidable in logarithmic space.
\end{theorem}

\begin{definition}[Labels of the bank]
\lean{Complexity.EPadBX}
\label{Complexity.EPadBX}
\leanok
The set of bank labels outside the validity block is a finite set comprising labels for counting
the input length, for scanning the outer pair, for copying and halving registers, for checking
the index word's tail, for decoding the index, for computing the seek limit, for the strided
seek, and the two answering labels $\mathsf{yes}$ and $\mathsf{no}$.
\end{definition}

\begin{definition}[Enumeration of the bank labels]
\lean{Complexity.epadBXAll}
\label{Complexity.epadBXAll}
\leanok
\uses{Complexity.EPadBX}
The enumeration of the bank's labels is an explicit finite list in which every bank label
outside the validity block appears, including each instance of the labels indexed by a bit or by
an element of $\{0,1,2,3\}$.
\end{definition}

\begin{definition}[Finiteness of the bank labels]
\lean{Complexity.epadBXFintype}
\label{Complexity.epadBXFintype}
\leanok
\uses{Complexity.EPadBX, Complexity.epadBXAll}
The set of bank labels outside the validity block is given the structure of a finite set, with
every label appearing in the explicit enumeration of the bank labels.
\end{definition}

\begin{definition}[Targets of the outer pair scan]
\lean{Complexity.epadZEnd}
\label{Complexity.epadZEnd}
\leanok
\uses{Complexity.EPadBX}
For $k \in \{0,1,2,3\}$, the target of advance $k$ of the bank's outer pair scan is the bank
label $\mathsf{ztest}_0$ if $k = 0$, $\mathsf{ztest}_1$ if $k = 1$, $\mathsf{zscan}$ if $k = 2$,
and $\mathsf{c1t}$ if $k = 3$.
\end{definition}

\begin{definition}[Bank register machine]
\lean{Complexity.epadBA}
\label{Complexity.epadBA}
\leanok
\uses{Complexity.EPadBX, Complexity.EPadVL, Complexity.LogProg.ARM, Complexity.epadValCode, Complexity.epadZEnd}
Let $C, d \in \mathbb{N}$ and $\mu \in \{0,1\}$. The bank is the abstract register machine with
$18$ registers and $2$ deciders (the input bit and input length deciders) that, on input $z$,
counts $|z|$ into register $9$, parses $z$ as a pair $\langle y, w\rangle$ (rejecting otherwise),
sets register $1$ to $2|y|$, rejects unless $w$ is empty or ends in $1$ and $|w|$ is at most the
bit length of $|y|$, decodes $\mathrm{val}(w)$ into register $13$, runs the validity block at
stride $2$ on $y$ (rejecting unless $y = \langle a, 1^{C(|a|+1)^d}\rangle$), and then seeks
position $\mathrm{val}(w)$ in steps of $4$: it rejects if $\mathrm{val}(w) \ge |a|$, and otherwise
accepts in mode $\mu = 0$ and answers letter $4\,\mathrm{val}(w)$ of $z$ in mode $\mu = 1$.
\end{definition}

\begin{theorem}[Correctness of the bank's final phase]
\lean{Complexity.epad_bank_post}
\label{Complexity.epad_bank_post}
\leanok
\uses{Complexity.LogProg.AHalt, Complexity.LogProg.AHalt.ret, Complexity.LogProg.AHalt.step, Complexity.LogProg.AReach, Complexity.LogProg.AReach.halt, Complexity.LogProg.AReach.refl, Complexity.LogProg.AReach.step, Complexity.LogProg.PreS, Complexity.LogProg.astep_call, Complexity.LogProg.astep_dec, Complexity.LogProg.epad_preS, Complexity.LogProg.indexedSeek_loop, Complexity.LogProg.registerCopy_loop, Complexity.LogProg.registerScale_loop, Complexity.LogProg.register_update_le, Complexity.LogProg.vword_whole₁, Complexity.epadBA, Complexity.epadIn_answer, Complexity.epadOr, Turing.pairEncode}
\difficulty{6}
Let $C, d \in \mathbb{N}$, $\mu \in \{0,1\}$, $z$ a binary string, and $B, n, i \in \mathbb{N}$
with $4n + 1 \le B$. Let $v$ be register contents with all registers at most $B$, $v(2) = n+1$,
$v(13) = i$ and $v(15) = v(16) = v(17) = 0$. Then the bank register machine in mode $\mu$, run on
input $z$ with the input deciders as oracle from label $\mathsf{q1t}$ with $v$, halts with answer
$0$ if $i \ge n$, and if $i < n$ with answer $1$ in mode $\mu = 0$ and letter $z_{4i}$ (with
$0$ beyond the end) in mode $\mu = 1$; moreover every configuration along the run satisfies the syntactic step preconditions (equality tests compare distinct registers and decider calls use pairwise distinct argument registers) and has all register values at most $B$.
\end{theorem}

\begin{theorem}[Correctness of the bank's first phase]
\lean{Complexity.epad_bank_pre}
\label{Complexity.epad_bank_pre}
\leanok
\uses{Complexity.EPadBX, Complexity.EPadVL, Complexity.LogProg.AReach, Complexity.LogProg.AReach.refl, Complexity.LogProg.AReach.step, Complexity.LogProg.AReach.trans, Complexity.LogProg.PreS, Complexity.LogProg.alignedPair_loop, Complexity.LogProg.astep_dec, Complexity.LogProg.astep_half, Complexity.LogProg.bitsVal, Complexity.LogProg.bitsVal_bits, Complexity.LogProg.bitsVal_lt_two_pow, Complexity.LogProg.boundedDecode_loop, Complexity.LogProg.canonicalTail_loop, Complexity.LogProg.epad_preS, Complexity.LogProg.inputLength_loop, Complexity.LogProg.registerCopy_loop, Complexity.LogProg.register_update_le, Complexity.LogProg.two_pow_bits_le, Complexity.epadBA, Complexity.epadIn_answer, Complexity.epadOr, Complexity.epadZEnd, Complexity.pairEncode_getD_snd, Turing.eq_pairEncode_of_pairDecode, Turing.length_pairEncode, Turing.pairDecode, Turing.pairEncode}
\difficulty{6}
Let $C, d \in \mathbb{N}$, $\mu \in \{0,1\}$, $z$ a binary string, and $B \ge 2|z| + 1$.
Run the bank register machine in mode $\mu$ on $z$, with the input deciders as oracle, from
label $\mathsf{lscan}$ with all registers $0$. If $z$ is not a self-delimiting pair, it reaches
$\mathsf{no}$. If $z = \langle y, w\rangle$ and it is not the case that $w$ is empty or ends in
$1$ and $|w|$ is at most the bit length of $|y|$, it reaches $\mathsf{no}$. If $z = \langle y,
w\rangle$ with $w$ satisfying both conditions, it reaches the validity block's $\mathsf{scan}$
label with register contents $v$ satisfying $v(0) = 0$, $v(1) = 2|y|$, $v(r) = 0$ for
$2 \le r \le 8$, $v(13) = \mathrm{val}(w)$ (the little-endian binary value of $w$) and
$v(15) = v(16) = v(17) = 0$. In each case the final registers are at most $B$ and every configuration along the run satisfies the syntactic step preconditions (equality tests compare distinct registers and decider calls use pairwise distinct argument registers) and has all register values at most $B$.
\end{theorem}

\begin{lemma}[Letter 4i of a nested pair]
\lean{Complexity.epad_getD_four}
\label{Complexity.epad_getD_four}
\leanok
\uses{Complexity.pairEncode_getD, Turing.length_pairEncode, Turing.pairEncode}
\difficulty{3}
Let $a, b, w$ be binary strings and $i < |a|$. Then letter $4i$ of the nested self-delimiting
pair $\langle\langle a, b\rangle, w\rangle$ equals the letter $a_i$.
\end{lemma}

\begin{theorem}[Bank languages are decidable in logspace]
\lean{Complexity.epad_Q_mem}
\label{Complexity.epad_Q_mem}
\leanok
\uses{Complexity.LOGSPACE, Complexity.LogProg.AHalt, Complexity.LogProg.AHalt.ret, Complexity.LogProg.AReach.halt, Complexity.LogProg.PreS, Complexity.LogProg.arm_decides_poly, Complexity.LogProg.bitsVal, Complexity.LogProg.epad_preS, Complexity.LogProg.indicator_eq_decide, Complexity.epadBA, Complexity.epadBXFintype, Complexity.epadIn, Complexity.epadIn_answer, Complexity.epadIn_mem, Complexity.epadOr, Complexity.epadQ, Complexity.epadS, Complexity.epadVLFintype, Complexity.epadValid, Complexity.epad_bank_post, Complexity.epad_bank_pre, Complexity.epad_bound, Complexity.epad_getD_four, Complexity.epad_mem_Q_iff, Complexity.epad_not_mem_Q, Complexity.epad_val_run, Complexity.pairEncode_getD, Complexity.pairFstD, Complexity.pairFstD_pairEncode, Turing.MultiTapeTM.indicator, Turing.eq_pairEncode_of_pairDecode, Turing.length_pairEncode, Turing.pairDecode, Turing.pairEncode}
\difficulty{6}
Let $C, d \in \mathbb{N}$ and $\mu \in \{0,1\}$. The bank language $Q_{C,d,\mu}$, of pairs
$\langle y, w\rangle$ with $w$ empty or ending in $1$, $|w|$ at most the bit length of $|y|$,
$y = \langle a, 1^{C(|a|+1)^d}\rangle$, and $\mathrm{val}(w) < |a|$ (mode $0$) or
$a_{\mathrm{val}(w)} = 1$ (mode $1$), belongs to $\mathsf{l}$, the class of languages decidable in
logarithmic space.
\end{theorem}

\begin{theorem}[Space-bounded simulation on the virtual input]
\lean{Complexity.epad_run}
\label{Complexity.epad_run}
\leanok
\uses{Complexity.LOGSPACE_eq, Complexity.LogProg.OIS.compile_decides, Complexity.LogProg.OIS.prog, Complexity.LogProg.OIS.start, Complexity.LogProg.bankK, Complexity.LogProg.bankStart, Complexity.LogProg.bankTM, Complexity.LogProg.compileFinTM, Complexity.SPACE, Complexity.epadQ, Complexity.epadValid, Complexity.epadVirt, Complexity.epadVirt_length_le, Complexity.epadVirt_of_not_valid, Complexity.epadVirt_of_valid, Complexity.epad_Q_mem, Complexity.epad_query_mem_iff, Complexity.logSpace, Complexity.logSpace_le_succ, Complexity.mem_SPACE_iff, Complexity.pairFstD_pairEncode, Turing.FinTM, Turing.FinTM.ComputesInSpace, Turing.FinTM.DecidesInSpace, Turing.FinTM.computesInSpace_iff, Turing.FinTM.decidesInSpace_iff, Turing.MultiTapeTM.indicator, Turing.MultiTapeTM.initCfg, Turing.MultiTapeTM.spaceUsed, Turing.length_bits_le_self, Turing.length_pairEncode, Turing.pairEncode}
\difficulty{6}
Let $S, T : \mathbb{N} \to \mathbb{N}$, where $\mathrm{SPACE}(f)$ denotes the class of languages
decided by a multi-tape Turing machine over the binary alphabet using at most $c\,f(n)$ work-tape
cells on inputs of length $n$, for some constant $c$. Let $L \in \mathrm{SPACE}(S)$, let $C, d \in \mathbb{N}$, and suppose that $T$ satisfies $S(n) \le T(2n + 2 + C(n+1)^d)$ for all $n$ and $m + 1 \le T(m)$ for all $m$. Then there
are a constant $K$ and a multi-tape Turing machine over the binary alphabet that, on every input
$y$, halts with output the single bit indicating whether $\mathrm{virt}(y) \in L$, having used at most $K\,T(|y|)$ work-tape cells, where
$\mathrm{virt}(y)$ is $a$ if $y = \langle a, 1^{C(|a|+1)^d}\rangle$ in the self-delimiting pair
encoding and the empty string otherwise.
\end{theorem}

\begin{theorem}[Padding transfers space bounds]
\lean{Complexity.mem_SPACE_padLang}
\label{Complexity.mem_SPACE_padLang}
\leanok
\uses{Complexity.LOGSPACE_eq, Complexity.LogProg.indicator_eq_decide, Complexity.SPACE, Complexity.epadV, Complexity.epadValid, Complexity.epadVirt, Complexity.epadVirt_of_valid, Complexity.epad_V_mem, Complexity.epad_mem_V, Complexity.epad_run, Complexity.logSpace, Complexity.logSpace_le_succ, Complexity.mem_SPACE_iff, Complexity.mem_padLang_iff, Complexity.padLang, Complexity.pairFstD, Complexity.pairFstD_pairEncode, Turing.FinTM.ComputesInSpace.computesFunInTime_and_forall_spaceUsed_le, Turing.FinTM.computesFunInTime_cond_spaceUsed, Turing.FinTM.computesFunInTime_const_spaceUsed, Turing.FinTM.computesFunInTime_iff, Turing.FinTM.computesInSpace_iff, Turing.FinTM.decidesInSpace_iff, Turing.MultiTapeTM.indicator}
\difficulty{5}
Let $S, T : \mathbb{N} \to \mathbb{N}$, where $\mathrm{SPACE}(f)$ denotes the class of languages
decided by a multi-tape Turing machine over the binary alphabet using at most $c\,f(n)$ work-tape
cells on inputs of length $n$, for some constant $c$. Let $L \in \mathrm{SPACE}(S)$, let $C, d \in \mathbb{N}$, and suppose that $T$ satisfies $S(n) \le T(2n + 2 + C(n+1)^d)$ for all $n$ and $m + 1 \le T(m)$ for all $m$. Then the
padded language $\{\langle x, 1^{C(|x|+1)^d}\rangle : x \in L\}$, with $\langle\cdot,\cdot\rangle$
the self-delimiting pair encoding (bits of $x$ doubled, then $01$, then the pad), lies in
$\mathrm{SPACE}(T)$.
\end{theorem}
